\chapter{Mathematical Physics}
\label{sec:prereq-physics}

This section provides prerequisite knowledge, definitions, theorems, and examples for the Mathematical Physics track.

\section{Classical Mechanics}

\textbf{Key Topics:} Lagrangian and Hamiltonian formalisms, central potentials, orbital stability, Binet's equation, Noether's theorem, and rotating frames.

\begin{definition}[Lagrangian and Hamiltonian]
The \emph{Lagrangian} $L$ of a classical system is defined as $L = T - V$, where $T$ is the kinetic energy and $V$ is the potential energy. The \emph{Hamiltonian} $H$ is obtained via a Legendre transformation: $H(q, p, t) = \sum_i p_i \dot{q}_i - L(q, \dot{q}, t)$, where $p_i = \frac{\partial L}{\partial \dot{q}_i}$ are the canonical momenta.
\end{definition}

\begin{theorem}[Noether's Theorem]
Every differentiable symmetry of the action of a physical system has a corresponding conservation law. For example, spatial translation invariance implies conservation of momentum, and time translation invariance implies conservation of energy.
\end{theorem}

\begin{definition}[Effective Potential and Stability]
In a central potential $V(r)$, the radial motion is governed by the effective potential $V_{\text{eff}}(r) = \frac{J^2}{2mr^2} + V(r)$, where $J = mr^2\dot{\phi}$ is the conserved angular momentum. A circular orbit exists at $r_0$ if $V_{\text{eff}}'(r_0) = 0$. The orbit is stable if $V_{\text{eff}}''(r_0) > 0$. Small oscillations around a stable orbit have frequency $\omega_{\text{osc}} = \sqrt{V_{\text{eff}}''(r_0)/m}$.
\end{definition}

\begin{definition}[Binet's Equation]
The orbit $u(\phi) = 1/r(\phi)$ of a particle in a central force $F(r)$ satisfies:
\[ \frac{d^2u}{d\phi^2} + u = -\frac{m}{J^2 u^2} F(1/u) \]
This equation is useful for determining the shape of the orbit without solving for time dependence.
\end{definition}

\begin{example}[Rotating Frame: Bead on a Hoop]
For a bead on a vertical hoop of radius $R$ rotating with angular velocity $\omega$, the Lagrangian in terms of the polar angle $\theta$ (from the bottom) is:
\[ L = \frac{1}{2}mR^2 \dot{\theta}^2 + \frac{1}{2}mR^2 \omega^2 \sin^2 \theta + mgR \cos \theta \]
The term $\frac{1}{2}mR^2 \omega^2 \sin^2 \theta$ represents the centrifugal potential. A pitchfork bifurcation occurs at $\omega = \sqrt{g/R}$, where the stable equilibrium shifts from the bottom ($\theta=0$) to a non-zero angle.
\end{example}

\section{Quantum Mechanics}

\textbf{Key Topics:} Harmonic oscillator, ladder operators, perturbation theory, coherent states, time reversal symmetry, and $su(2)$ representations.

\begin{definition}[Ladder Operators and Coherent States]
For a 1D harmonic oscillator, $\hat{a} = \sqrt{\frac{m\omega}{2\hbar}} \hat{x} + \frac{i}{\sqrt{2m\hbar\omega}} \hat{p}$. A \emph{coherent state} $|\alpha\rangle$ is an eigenstate of $\hat{a}$ with $\hat{a}|\alpha\rangle = \alpha|\alpha\rangle$. It can be expressed as $|\alpha\rangle = e^{-|\alpha|^2/2} \sum_{n=0}^\infty \frac{\alpha^n}{\sqrt{n!}} |n\rangle$.
\end{definition}

\begin{definition}[Anti-unitary Operators and Time Reversal]
An operator $\Theta$ is \emph{anti-unitary} if it is anti-linear ($\Theta(c|\psi\rangle) = c^* \Theta|\psi\rangle$) and satisfies $\langle \Theta \psi | \Theta \phi \rangle = \langle \phi | \psi \rangle$. The time reversal operator $\Theta$ is anti-unitary. If $[H, \Theta] = 0$ and an energy eigenstate is non-degenerate, its wavefunction can be chosen to be real.
\end{definition}

\begin{theorem}[Schwinger Boson Representation of $su(2)$]
Given two independent oscillators with ladder operators $a_i, a_i^\dagger$ ($i=1,2$), the operators
\[ J_+ = a_1^\dagger a_2, \quad J_- = a_2^\dagger a_1, \quad J_z = \frac{1}{2}(a_1^\dagger a_1 - a_2^\dagger a_2) \]
satisfy the $su(2)$ commutation relations: $[J_z, J_\pm] = \pm J_\pm$ and $[J_+, J_-] = 2J_z$. For a fixed total number of quanta $n = n_1 + n_2$, these states form an irreducible representation with spin $j = n/2$.
\end{theorem}

\begin{example}[Degenerate Perturbation Theory]
Consider a 2D isotropic harmonic oscillator $H_0 = \hbar\omega(n_x + n_y + 1)$ with perturbation $V = \alpha m \omega^2 x y$. For the first excited states $\{|1,0\rangle, |0,1\rangle\}$, the perturbation matrix is $V = \frac{1}{2}\alpha \hbar \omega \begin{pmatrix} 0 & 1 \\ 1 & 0 \end{pmatrix}$, leading to energy shifts $\Delta E = \pm \frac{1}{2}\alpha \hbar \omega$.
\end{example}

\section{Electromagnetism}

\textbf{Key Topics:} Maxwell equations, wave propagation in conductors, gauge transformations, retarded potentials, and the Lorentz gauge.

\begin{definition}[Maxwell Equations]
In vacuum (Gaussian units), Maxwell's equations are:
\begin{enumerate}
    \item $\nabla \cdot \mathbf{E} = 0$, $\nabla \cdot \mathbf{B} = 0$
    \item $\nabla \times \mathbf{E} = -\frac{1}{c}\frac{\partial \mathbf{B}}{\partial t}$, $\nabla \times \mathbf{B} = \frac{1}{c}\frac{\partial \mathbf{E}}{\partial t}$
\end{enumerate}
\end{definition}

\begin{definition}[Electromagnetic Potentials and Gauge Invariance]
Physical fields are derived from potentials $\mathbf{E} = -\nabla \phi - \frac{1}{c} \frac{\partial \mathbf{A}}{\partial t}$ and $\mathbf{B} = \nabla \times \mathbf{A}$. A \emph{gauge transformation} $\mathbf{A} \to \mathbf{A} + \nabla \Lambda$, $\phi \to \phi - \frac{1}{c} \frac{\partial \Lambda}{\partial t}$ leaves $\mathbf{E}$ and $\mathbf{B}$ invariant. The \emph{Lorentz gauge} $\nabla \cdot \mathbf{A} + \frac{1}{c} \frac{\partial \phi}{\partial t} = 0$ simplifies the wave equations for potentials.
\end{definition}

\begin{theorem}[Retarded Potentials]
The solutions to the inhomogeneous wave equations in the Lorentz gauge are:
\[ \phi(\mathbf{r}, t) = \int \frac{\rho(\mathbf{r}', t_r)}{|\mathbf{r} - \mathbf{r}'|} d^3 r', \quad \mathbf{A}(\mathbf{r}, t) = \frac{1}{c} \int \frac{\mathbf{J}(\mathbf{r}', t_r)}{|\mathbf{r} - \mathbf{r}'|} d^3 r' \]
where $t_r = t - |\mathbf{r} - \mathbf{r}'|/c$ is the retarded time.
\end{theorem}

\begin{theorem}[Propagation in Conductors]
In a medium with high conductivity $\sigma \gg \omega$, the wave vector $k$ is complex: $k \approx \frac{1+i}{\sqrt{2}} \frac{\sqrt{4\pi \sigma \omega}}{c}$. The real part describes propagation, and the imaginary part describes attenuation with a characteristic \emph{skin depth} $\delta = c/\sqrt{2\pi \sigma \omega}$.
\end{theorem}

\begin{example}[Finite Beams and Longitudinal Components]
For a wave packet with transverse profile $E_0(x,y)$ propagating in $z$, the transversal field $\mathbf{E}_\perp = E_0(x,y) e^{i(kz-\omega t)} \frac{\hat{\mathbf{x}}+i\hat{\mathbf{y}}}{\sqrt{2}}$ is not a complete solution. Gauss's law $\nabla \cdot \mathbf{E} = 0$ requires a longitudinal component $E_z$:
\[ E_z = \frac{i}{k\sqrt{2}}(\partial_x E_0 + i \partial_y E_0) e^{i(kz-\omega t)} \]
This component contributes to the beam's orbital angular momentum.
\end{example}

\begin{definition}[Energy and Momentum]
The time-averaged energy density is $\langle u \rangle = \frac{1}{16\pi}(|\mathbf{E}|^2 + |\mathbf{B}|^2)$ and the momentum density is $\langle \mathbf{g} \rangle = \frac{1}{8\pi c} \text{Re}(\mathbf{E} \times \mathbf{B}^*)$. For a circularly polarized Gaussian beam, the ratio of total angular momentum $\langle L_z \rangle$ to total energy $\langle U \rangle$ is $1/\omega$.
\end{definition}

\section{Statistical Mechanics and General Relativity}

\textbf{Key Topics:} Partition function, Fermi gas, Bose-Einstein statistics, density of states, Schwarzschild metric, Einstein field equations, and conformal coupling.

\begin{definition}[Partition Function and Distribution Functions]
The canonical partition function is $Z = \sum_s e^{-\beta E_s}$, with $\beta = 1/k_B T$. The average occupation number for particles in state $i$ with energy $E_i$ and chemical potential $\mu$ is:
\begin{enumerate}
    \item $n_B(E_i) = \frac{1}{e^{\beta(E_i-\mu)} - 1}$ (Bose-Einstein statistics)
    \item $n_F(E_i) = \frac{1}{e^{\beta(E_i-\mu)} + 1}$ (Fermi-Dirac statistics)
\end{enumerate}
\end{definition}

\begin{definition}[Density of States]
The density of states $g(E)$ represents the number of states per unit energy. For a 1D system of length $L$, $g(E) = \frac{L}{2\pi} \frac{dk}{dE}$. For massless particles ($E = \hbar c k$), $g(E) = \frac{L}{2\pi \hbar c}$. For massive particles ($E = \hbar^2 k^2 / 2m$), $g(E) = \frac{L}{2\pi} \sqrt{\frac{2m}{\hbar^2 E}}$.
\end{definition}

\begin{definition}[Einstein Field Equations and Metrics]
The Einstein field equations relate spacetime geometry to energy-momentum: $R_{\mu\nu} - \frac{1}{2} R g_{\mu\nu} = 8\pi G T_{\mu\nu}$. In vacuum, $R_{\mu\nu} = 0$. The Schwarzschild metric is the unique spherically symmetric vacuum solution:
\[ ds^2 = -\left(1 - \frac{2M}{r}\right) dt^2 + \left(1 - \frac{2M}{r}\right)^{-1} dr^2 + r^2 d\Omega^2 \]
A Killing vector $K^\mu$ generates a symmetry of the metric, $\nabla_{(\mu} K_{\nu)} = 0$.
\end{definition}

\begin{definition}[Conformal Transformations and Coupling]
A \emph{conformal transformation} is a local rescaling of the metric $g_{\mu\nu} \to \Omega^2(x) g_{\mu\nu}$. A scalar field action $\int d^4x \sqrt{-g} [\frac{1}{2}(\partial \phi)^2 + \frac{1}{2} \xi R \phi^2]$ is conformally invariant if $\xi = 1/6$ in 4D. In de Sitter space (constant $R$), this coupling introduces an effective mass $m^2 = \xi R$.
\end{definition}

\section{Quantum Field Theory}

\textbf{Key Topics:} Feynman diagrams, renormalization, scalar field theories, and conformal invariance.

\begin{definition}[Feynman Diagrams]
Feynman diagrams graphically represent the perturbative expansion of scattering amplitudes in QFT. Propagators represent free particles, while vertices correspond to interactions between them.
\end{definition}

\begin{theorem}[Dimensional Regularization]
To handle divergent loop integrals in QFT, dimensional regularization evaluates integrals in $d=4-\epsilon$ dimensions. Divergences appear as $1/\epsilon$ poles, which can be systematically subtracted (e.g., in the $\overline{\text{MS}}$ scheme) via renormalization to yield finite physical predictions.
\end{theorem}

\begin{example}[One-Loop Self-Energy in $\phi^4$ Theory]
In massless scalar $\phi^4$ theory, the interaction term is $-\frac{\lambda}{4!} \phi^4$. The one-loop correction to the propagator is quadratically divergent in 4 dimensions. Renormalization introduces a mass counterterm to absorb this divergence, allowing for a well-defined renormalized mass.
\end{example}